%! Tex main = MarquezSalazarBrandon-PropuestaDeTesis.tex
\subsection{Justificación}

La exploración de configuraciones de vacío en teoría de cuerdas
constituye un problema de optimización en espacios de alta
dimensionalidad, donde las funciones de energía potencial
presentan múltiples óptimos locales y regiones caóticas que
dificultan el uso de métodos analíticos tradicionales o
gradiente-dependientes \cite{Lanza_2025, De_Luca_2025}. En este
contexto, las metaheurísticas ofrecen una alternativa natural
para abordar la búsqueda sistemática de puntos estables.

En optimización computacional, el teorema de
\foreignlanguage{english}{\textit{No Free Lunch}} (NFL) establece que no
existe un algoritmo universalmente superior para todos los
problemas
\cite{Wolpert_1997}. El desempeño relativo de los métodos
depende fuertemente de las características específicas del
problema abordado. Para el caso de la optimización de funciones
potenciales derivadas de la K-teoría, no existe garantía
\textit{a priori} de que un método particular sea la estrategia más
adecuada en términos de eficiencia, precisión o capacidad de
exploración.

Por tanto, resulta pertinente explorar un conjunto amplio de
metaheurísticas (algoritmos genéticos, optimización por
enjambre de partículas, algoritmos de estimación de
distribución, colonia de hormigas, entre otros) para determinar
cuáles se adaptan mejor a la geometría del problema.
La importancia de este proyecto radica en lograr una mejora en
el método de búsqueda: al implementar y evaluar múltiples
enfoques, se busca incrementar la cantidad de vacíos estables
identificados respecto a trabajos previos como \cite{Damian_2022},
ampliando el catálogo de configuraciones conocidas y
contribuyendo a una comprensión más completa del paisaje de
soluciones.
